\documentclass[11pt]{article}
\usepackage[a4paper,margin=27mm]{geometry}
\usepackage{amsmath,amssymb,amsthm,mathtools}
\usepackage{booktabs,longtable,array,enumitem,microtype}
\usepackage{xcolor}
\usepackage[T1]{fontenc}
\usepackage{lmodern}
\usepackage[hidelinks]{hyperref}

\definecolor{proved}{RGB}{0,105,75}
\definecolor{conditional}{RGB}{160,92,0}
\definecolor{open}{RGB}{165,25,35}
\newcommand{\Pstatus}{\textcolor{proved}{\textsf{PROVED}}}
\newcommand{\Cstatus}{\textcolor{conditional}{\textsf{CONDITIONAL}}}
\newcommand{\Ostatus}{\textcolor{open}{\textsf{OPEN}}}
\newcommand{\Fstatus}{\textcolor{open}{\textsf{FALSIFIED}}}
\newcommand{\TV}{\operatorname{TV}}
\newcommand{\Cov}{\operatorname{Cov}}
\newcommand{\diam}{\operatorname{diam}}
\newcommand{\e}{\mathrm e}
\newcommand{\cA}{\mathcal A}
\newcommand{\cG}{\mathcal G}
\newcommand{\cH}{\mathcal H}

\newtheorem{definition}{Definition}[section]
\newtheorem{theorem}[definition]{Theorem}
\newtheorem{proposition}[definition]{Proposition}
\newtheorem{lemma}[definition]{Lemma}
\newtheorem{corollary}[definition]{Corollary}
\newtheorem{obligation}[definition]{Frontier obligation}
\theoremstyle{remark}
\newtheorem{remark}[definition]{Remark}

\title{Heat-Kernel Boundary Sensitivity and a Two-Tier\\
Infrared Basin for Lattice Yang--Mills}
\author{Research dossier}
\date{20 July 2026}

\begin{document}
\maketitle

\begin{abstract}
We implement the first theorem-level step of a representation-aware programme
for the Yang--Mills mass gap.  The heat-kernel plaquette action is selected as
a primary regulator because gauge invariance, positive character coefficients,
and convolution structure are explicit.  We prove that a naive worst-case
total-variation Dobrushin basin cannot be uniform in the small heat-time limit:
on every nontrivial compact group, a heat kernel and any fixed nontrivial
translate become asymptotically mutually singular.  This is a constructive
obstruction, not a no-go theorem for clustering.  It shows that continuum basin
entry must combine scale-sensitive good-field couplings with quantitative
control of exceptional and topological blocks.  We define such a two-tier
basin, prove its abstract basin-to-clustering and stability theorems, give an
executable $U(1)$ calibration, and isolate two explicit four-dimensional entry
inequalities.  Their verification for pure Yang--Mills remains open.
\end{abstract}

\tableofcontents

\section{Target and result status}

The ultimate target is the construction of nontrivial four-dimensional pure
Yang--Mills theory with a Hamiltonian gap above the vacuum for every compact
simple gauge group \cite{JaffeWitten}.  The present paper addresses one earlier
interface: a local, regulator-level criterion that could generate exponential
infrared rigidity after renormalization.

The results are:
\begin{enumerate}[label=(\roman*)]
\item a finite-lattice heat-kernel regulator and positivity certificate;
\item a compact-group obstruction theorem for worst-case total-variation
      influence;
\item a two-tier good-field/defect comparison theorem;
\item explicit replacement basin-entry inequalities; and
\item tested $U(1)$ and $SU(2)$ heat-kernel calibrations, an exact $U(1)$
      convolution-blocking experiment, and a curvature-matched Wilson
      comparison.
\end{enumerate}
Items (ii)--(iii) are rigorous abstract results.  The proof that the
four-dimensional Yang--Mills RG trajectory satisfies the replacement
inequalities is not supplied.  Consequently no continuum construction or
Clay-level mass gap is claimed.

\section{Selection of the heat-kernel regulator}

Let $G$ be a compact connected Lie group with normalized Haar measure $dg$ and
bi-invariant Laplacian $\Delta_G$.  Let $p_t(g)$ be the heat kernel at time
$t>0$, normalized by
\[
 \partial_t p_t=\frac12\Delta_Gp_t,
 \qquad \int_Gp_t(g)\,dg=1.
\]
It is central, satisfies $p_t(g^{-1})=p_t(g)$, and has the character expansion
\begin{equation}\label{eq:character}
 p_t(g)=\sum_{\lambda\in\widehat G}
 d_\lambda\e^{-tC_2(\lambda)/2}\chi_\lambda(g),
\end{equation}
with positive coefficients.  This action was introduced for lattice gauge
theory by Menotti and Onofri \cite{MenottiOnofri}.

On a finite oriented lattice $\Lambda=(V,E,P)$, assign $U_e\in G$ to each
oriented edge with $U_{e^{-1}}=U_e^{-1}$.  For an oriented plaquette $p$, let
$U_p$ be the ordered boundary holonomy.  Define
\begin{equation}\label{eq:measure}
 d\mu_{t,\Lambda}(U)
 =Z_{t,\Lambda}^{-1}\prod_{p\in P}p_t(U_p)
   \prod_{e\in E^+}dU_e.
\end{equation}

\begin{proposition}[Finite gauge invariance \Pstatus]
The measure \eqref{eq:measure} is invariant under
$U_{xy}\mapsto g_xU_{xy}g_y^{-1}$ for arbitrary $(g_x)_{x\in V}\in G^V$.
\end{proposition}

\begin{proof}
Plaquette holonomy transforms by conjugation at its base point.  The heat
kernel is central, and Haar measure is left- and right-invariant.
\end{proof}

\begin{proposition}[Finite reflection-positive factor \Pstatus]\label{prop:rp}
Consider a lattice reflection that splits every plaquette crossing the
reflection plane into a negative and a positive half-holonomy.  Each crossing
heat-kernel factor has a positive sum-of-squares decomposition.  Consequently,
with reflection-compatible boundary conditions, the standard finite-lattice
reflection-positivity form is nonnegative.
\end{proposition}

\begin{proof}
Write a crossing holonomy as $a^{-1}b$, where reflection exchanges the half
variables.  By \eqref{eq:character},
\[
 p_t(a^{-1}b)=\sum_{\lambda,i,j}
 d_\lambda\e^{-tC_2(\lambda)/2}
 \overline{\lambda(a)_{ij}}\lambda(b)_{ij}.
\]
The coefficient of every matrix element is nonnegative.  Products of crossing
factors retain a Gram decomposition; plaquettes away from the plane occur in
reflected pairs.  Integrating a positive-time cylinder function against its
reflection therefore gives a sum of squared norms.
\end{proof}

This proposition is a finite-regulator certificate.  RG blocking must either
preserve the factorization or carry an independent recovery certificate.

\subsection{Why this regulator is selected}

\begin{center}
\small
\begin{tabular}{p{.22\textwidth}p{.31\textwidth}p{.34\textwidth}}
\toprule
Feature & Advantage & Remaining obligation\\
\midrule
central heat kernel & exact gauge invariance & global orbit and topology control\\
positive characters & transparent finite positivity & preservation or recovery under RG\\
semigroup convolution & natural block comparison in two dimensions and toy sectors & no exact four-dimensional RG closure is assumed\\
small-time local Gaussianity & connects to Lie-algebra weak fields & worst-case boundary sensitivity degenerates\\
strictly positive density & finite conditional kernels are well defined & uniform physical-scale inequalities remain open\\
\bottomrule
\end{tabular}
\end{center}

The Wilson regulator is retained as an adversarial comparison, not as a
foundation.  A universality theorem between the two remains part of the final
continuum programme.

\section{The original basin and its hidden quantifier}

For a Gibbs specification $(\gamma_i)$ on blocks $i\in I$, the classical
worst-case total-variation coefficient is schematically
\begin{equation}\label{eq:naive}
 c_{ij}=\sup_{x=z\text{ off }j}
 \TV\!\left(\gamma_i(\cdot\mid x),
             \gamma_i(\cdot\mid z)\right).
\end{equation}
A uniform bound $\sup_i\sum_jc_{ij}<1$ implies uniqueness and comparison
estimates in its domain of applicability.  Generalized block comparison
theorems can extend that domain \cite{RebeschiniVanHandel}.

The critical word in \eqref{eq:naive} is ``sup''.  In a weak-field limit the
conditional law can become extremely narrow.  Boundary data that move its
center by a fixed group element then produce nearly disjoint conditionals even
when typical neighboring configurations are much more coherent.

\section{Compact-group heat-kernel obstruction}

We use the convention
\[
 \TV(\mu,\nu)=\sup_{A}|\mu(A)-\nu(A)|
              =\frac12\|\mu-\nu\|_{L^1}
\]
when densities exist.

\begin{lemma}[Disjoint concentration \Pstatus]\label{lem:disjoint}
Let $X$ be a Hausdorff space and let $\mu_t\Rightarrow\delta_x$ and
$\nu_t\Rightarrow\delta_y$ weakly as $t\downarrow0$, with $x\ne y$.
Then $\TV(\mu_t,\nu_t)\to1$.
\end{lemma}

\begin{proof}
Choose disjoint open neighborhoods $U\ni x$ and $V\ni y$ whose closures remain
disjoint.  By the Portmanteau theorem, after shrinking if necessary,
$\mu_t(U)\to1$ and $\nu_t(U)\to0$.  Thus
$\TV(\mu_t,\nu_t)\ge|\mu_t(U)-\nu_t(U)|\to1$, while total variation is at most
one.
\end{proof}

\begin{theorem}[Translated heat-kernel obstruction \Pstatus]\label{thm:tv}
Let $G$ be a nontrivial compact Lie group, let
$\mu_t(dg)=p_t(g)dg$, and fix $h\ne e$.  If $L_h\mu_t$ denotes its left
translate, then
\[
 \lim_{t\downarrow0}\TV(\mu_t,L_h\mu_t)=1.
\]
\end{theorem}

\begin{proof}
The heat kernels form an approximate identity, hence
$\mu_t\Rightarrow\delta_e$.  Translation is continuous, so
$L_h\mu_t\Rightarrow\delta_h$.  Apply Lemma~\ref{lem:disjoint}.
\end{proof}

\begin{corollary}[Failure of uniform worst-case contraction \Fstatus]
Consider the one-plaquette conditional family
\[
 \gamma_t^\eta(dx)=p_t(x\eta^{-1})\,dx,
 \qquad\eta\in G.
\]
Then
\[
 c(t):=\sup_{\eta,\eta'}\TV(\gamma_t^\eta,\gamma_t^{\eta'})
 \longrightarrow1
 \quad(t\downarrow0).
\]
In particular, there is no $q<1$ such that $c(t)\le q$ for every sufficiently
small $t$.
\end{corollary}

\begin{proof}
Choose $\eta'=h\eta$ with fixed $h\ne e$.  The two conditional laws are
translates covered by Theorem~\ref{thm:tv}.
\end{proof}

\begin{remark}
The theorem does not prove absence of a mass gap, uniqueness, or clustering.
It falsifies a proposed proof resource: a cutoff-uniform worst-case
single-plaquette total-variation coefficient.  Larger blocks, quotient
metrics, typical-set restrictions, and defect expansions remain available.
\end{remark}

\section{Executable circle calibration}

For $G=U(1)$ with angle $\theta$, the heat kernel is
\begin{equation}\label{eq:circle}
 p_t(\theta)=\frac1{2\pi}
 \left(1+2\sum_{n\ge1}\e^{-n^2t/2}\cos(n\theta)\right).
\end{equation}
The supplied program truncates this positive character series at a certified
coefficient tolerance and applies periodic trapezoidal quadrature.

\begin{center}
\begin{tabular}{rrrr}
\toprule
$t$ & $\TV(p_t,p_t(\cdot-\pi))$ &
$\TV(p_t,p_t(\cdot-0.5\sqrt t))$ & Gaussian limit\\
\midrule
1.00 & 0.767544967 & 0.194274584 & 0.197412651\\
0.50 & 0.947357851 & 0.197387803 & 0.197412651\\
0.20 & 0.999111867 & 0.197412653 & 0.197412651\\
0.10 & 0.999998642 & 0.197412651 & 0.197412651\\
0.05 & 1.000000000 & 0.197412642 & 0.197412651\\
0.02 & 1.000000000 & 0.197412640 & 0.197412651\\
\bottomrule
\end{tabular}
\end{center}

The final column is
\[
 \TV(N(0,1),N(0.5,1))
 =\operatorname{erf}\left(\frac{0.5}{2\sqrt2}\right).
\]
The fixed-shift computation illustrates Theorem~\ref{thm:tv}.  The
$\sqrt t$ shift shows that merely rescaling the metric does not automatically
produce a small coefficient: one obtains a nontrivial tangent-Gaussian
influence.  Actual contraction must come from block dynamics or a smaller
response of the conditional center to boundary perturbations.

Seven unit tests check positivity at the calibration grid, zero-shift identity,
translation symmetry, fixed-shift degeneration, the tangent-Gaussian limit,
output schema, and fail-closed parameter validation.

\section{Advanced regulator calibrations}

\subsection{$SU(2)$ character calibration}

For $SU(2)$, write the conjugacy angle as $\theta\in[0,\pi]$.  With the
Casimir convention used by the executable artifact,
\begin{equation}\label{eq:su2}
 p_t^{SU(2)}(\theta)
 =\sum_{n\ge0}(n+1)\e^{-tn(n+2)/2}
   \frac{\sin((n+1)\theta)}{\sin\theta},
 \qquad
 dg_{\rm class}=\frac2\pi\sin^2\theta\,d\theta.
\end{equation}
Translation by the nontrivial central element $-I$ sends
$\theta\mapsto\pi-\theta$.  Midpoint Haar quadrature gives
\begin{center}
\begin{tabular}{rr}
\toprule
$t$ & $\TV(p_t^{SU(2)}(\theta),p_t^{SU(2)}(\pi-\theta))$\\
\midrule
0.50 & 0.770181734\\
0.20 & 0.991849661\\
0.10 & 0.999976815\\
0.05 & 1.000000000\\
\bottomrule
\end{tabular}
\end{center}
This completes the simple-group numerical calibration of
Theorem~\ref{thm:tv}; the theorem itself does not depend on the truncation.

\subsection{Exact $U(1)$ convolution blocking}

The circle heat kernel obeys $p_t^{*m}=p_{mt}$.  Therefore a one-dimensional
chain or exactly convolutional plaquette block provides a controlled diagnostic
of how fixed-shift sensitivity changes under coarse graining:
\begin{center}
\begin{tabular}{rrr}
\toprule
Steps $m$ & coarse time $mt$ for $t=0.02$ & antipodal TV\\
\midrule
1 & 0.02 & 1.000000000\\
2 & 0.04 & 1.000000000\\
5 & 0.10 & 0.999998642\\
10 & 0.20 & 0.999111867\\
20 & 0.40 & 0.973991076\\
50 & 1.00 & 0.767544967\\
\bottomrule
\end{tabular}
\end{center}
Blocking reduces the fixed-shift coefficient only after the accumulated heat
time becomes appreciable.  Four-dimensional gauge blocking is not an exact
plaquette convolution, so this is a mechanism calibration rather than an RG
theorem.

\subsection{Defect-tail and regulator-comparison diagnostics}

Define a circle good field by $|\theta|\le2\sqrt t$.  Let $R_{\rm tail}$ be its
complement probability and let $K_{0.25}$ be the TV response to a
$0.25\sqrt t$ boundary shift.  The table also compares the heat kernel with the
curvature-matched Wilson density $Z^{-1}\exp(t^{-1}\cos\theta)$.
\begin{center}
\small
\begin{tabular}{rrrrr}
\toprule
$t$ & HK/Wilson TV & $R_{\rm tail}$ & $K_{0.25}$ & $8(K_{0.25}+R_{\rm tail})$\\
\midrule
0.50 & 0.063594784 & 0.045482439 & 0.099465621 & 1.159584482\\
0.20 & 0.019917856 & 0.045528368 & 0.099476451 & 1.160038549\\
0.10 & 0.009231571 & 0.045524826 & 0.099476455 & 1.160010248\\
0.05 & 0.004476830 & 0.045528348 & 0.099476460 & 1.160038463\\
\bottomrule
\end{tabular}
\end{center}
The decreasing second column is a one-plaquette comparison, not a universality
proof.  The final column is deliberately labeled an \emph{optimistic proxy}:
it treats a raw defect probability as though it were the required Lipschitz
defect influence and uses degree eight as a four-dimensional nearest-neighbor
stress test.  Its value above one shows that this toy choice of good radius and
local shift does not enter even the replacement basin.  A successful argument
must gain contraction through a genuine block quotient coupling, a sharper
defect-disagreement estimate, or both.

Together the basic and advanced suites contain fourteen passing tests.

\section{Two-tier good-field/defect basin}

The obstruction suggests splitting two effects that the worst-case supremum
mixes together: smooth response inside a regular orbit chart and propagation
through exceptional or topological blocks.

Let $I$ be a finite or countable block graph with metric $d$.  Let each block
state space $(S_i,\rho_i)$ have diameter at most one; for gauge systems it is a
quotient or pseudometric on physical block data.  Let $\gamma_i(\cdot\mid x)$
be local conditional kernels.  Choose gauge-invariant good events
$\mathsf G_i$.

\begin{definition}[Two-tier local coupling certificate]
For boundary configurations $x,z$ differing only at block $j$, a coupling
$Q_{ij}^{x,z}$ of the two $i$-conditionals has a certificate $(K_{ij},R_{ij})$
if
\begin{align}
 \mathbb E_Q[\rho_i(X_i,Z_i)\,;
       X_i,Z_i\in\mathsf G_i]
 &\le K_{ij}\rho_j(x_j,z_j),\label{eq:good}\\
 \mathbb E_Q[\rho_i(X_i,Z_i)\,;
       \{X_i,Z_i\}\not\subset\mathsf G_i]
 &\le R_{ij}\rho_j(x_j,z_j).
 \label{eq:bad}
\end{align}
Here $K$ measures regular transport and $R$ dominates defect activation,
topological mismatch, and chart failure.  Both are local quantities; neither is
defined by long-distance correlation decay.
\end{definition}

The estimates imply a Wasserstein influence matrix $M\le K+R$.  In a concrete
construction, \eqref{eq:bad} can follow from a disagreement-polymer estimate
rather than pointwise convexity.

\begin{definition}[Two-tier basin]
For $\kappa>0$ and $q<1$, a specification lies in
$\mathfrak B_{\rm 2t}(\kappa,q)$ when it has a two-tier certificate and
\begin{equation}\label{eq:twotier}
 \|K+R\|_\kappa
 :=\sup_i\sum_j\e^{\kappa d(i,j)}(K_{ij}+R_{ij})\le q<1,
\end{equation}
uniformly in total volume.
\end{definition}

\begin{lemma}[Certificate combination \Pstatus]
The full coupling satisfies
\[
 W_{1,\rho_i}(\gamma_i(\cdot\mid x),\gamma_i(\cdot\mid z))
 \le (K_{ij}+R_{ij})\rho_j(x_j,z_j).
\]
\end{lemma}

\begin{proof}
Add \eqref{eq:good} and \eqref{eq:bad}; together their events partition the
coupling space.  The Wasserstein distance is bounded by the cost of this
coupling.
\end{proof}

\begin{theorem}[Two-tier basin-to-clustering \Cstatus]\label{thm:twotier}
Assume the local specification admits the Wasserstein comparison theorem for
the seminorms
\[
 \delta_iF=\sup_{x=z\text{ off }i}
 \frac{|F(x)-F(z)|}{\rho_i(x_i,z_i)}.
\]
If it lies in $\mathfrak B_{\rm 2t}(\kappa,q)$, then for local $F,G$,
\begin{equation}\label{eq:cluster}
 |\Cov(F,G)|
 \le \frac{C_0}{1-q}\e^{-\kappa d(S_F,S_G)}
 \left(\sum_i\delta_iF\right)
 \left(\sum_j\delta_jG\right),
\end{equation}
where $C_0$ is the normalization constant in the chosen comparison theorem.
\end{theorem}

\begin{proof}
Let $M=K+R$.  The weighted row norm is submultiplicative, so
$D=(I-M)^{-1}=\sum_{n\ge0}M^n$ and $\|D\|_\kappa\le(1-q)^{-1}$.  Hence
\[
 D_{ij}\le(1-q)^{-1}\e^{-\kappa d(i,j)}.
\]
The Wasserstein comparison theorem bounds covariance by
$C_0\sum_{ij}\delta_iF D_{ij}\delta_jG$.  Nonzero seminorms lie in
$S_F$ and $S_G$, which yields \eqref{eq:cluster}.
\end{proof}

The comparison-theorem hypothesis is explicit because Gauss constraints can
prevent a naive product-site theorem from applying directly.  One must either
formulate the blocks on physical quotient data or prove the constrained
version required by the chosen regulator.

\begin{proposition}[Openness and perturbative stability \Pstatus]
Suppose $(K,R)$ satisfies \eqref{eq:twotier} with margin $1-q$ and another
effective theory has certificates $(K',R')$ with
\[
 \|(K'-K)+(R'-R)\|_\kappa\le\varepsilon<1-q.
\]
Then it lies in $\mathfrak B_{\rm 2t}(\kappa,q+\varepsilon)$.
\end{proposition}

\begin{proof}
Apply the triangle inequality for the weighted row norm.
\end{proof}

\begin{corollary}[Physical decay rate \Cstatus]
If the basin is reached at physical block length $\ell_*$ and observable
pullback errors are summable, the corresponding original observables have
correlation rate at least $\kappa/\ell_*$.  If the controlled separation is
Euclidean time and OS reconstruction plus observable density hold, the
positive spectral argument yields the same lower bound for the Hamiltonian
gap.
\end{corollary}

\section{Gauge-invariant good blocks and defect data}

For a block $B$, a regulator-level good event can be defined without a global
gauge slice by
\begin{equation}\label{eq:goodblock}
 \mathsf G_B(\rho)
 =\left\{U:\sum_{p\subset B}
 d_G(U_p,e)^2\le\rho^2|B|\right\},
\end{equation}
supplemented by physical bounds on boundary holonomy and an atlas label when a
single orbit chart is insufficient.  Plaquette conjugation invariance makes
\eqref{eq:goodblock} gauge invariant.

The defect tier records at least:
\begin{itemize}
\item high-curvature blocks violating \eqref{eq:goodblock};
\item failure of the chosen orbit chart or large transition functions between
      charts;
\item nontrivial center flux, global holonomy, or bundle-sector mismatch;
\item instanton-like concentration relevant at the selected scale; and
\item RG remainder events whose reconstruction error exceeds its budget.
\end{itemize}

A volume-small probability is insufficient.  What is required is a connected
polymer estimate whose entropy is dominated uniformly in volume and across the
number of RG steps used to reach $\ell_*$.

\section{Revised basin-entry theorem target}

Let $a$ be the microscopic spacing, $b>1$ the block factor, and
$\ell_k=b^ka$.  Let $K^{(k)}$ and $R^{(k)}$ be the regular and defect matrices
for the effective heat-kernel representation after $k$ controlled steps.

\begin{obligation}[Good-field contraction $\mathrm{BE_G}$ \Ostatus]
Prove that for some finite physical $\ell_*$ and $\kappa>0$,
\begin{equation}\label{eq:BEG}
 \sup_B\sum_C\e^{\kappa d(B,C)}K^{(k_*)}_{BC}le q_G,
\end{equation}
uniformly in cutoff and volume.  The matrix must be obtained from couplings on
gauge-orbit block data and must include all generated quasilocal interactions.
\end{obligation}

\begin{obligation}[Defect subcriticality $\mathrm{BE_D}$ \Ostatus]
Construct a defect-disagreement matrix or connected-polymer domination with
\begin{equation}\label{eq:BED}
 \sup_B\sum_C\e^{\kappa d(B,C)}R^{(k_*)}_{BC}le q_D,
 \qquad q_G+q_D<1,
\end{equation}
uniformly in cutoff and volume, including global orbit and topological data.
\end{obligation}

\begin{theorem}[Revised conditional entry theorem \Cstatus]
If $\mathrm{BE_G}$ and $\mathrm{BE_D}$ hold, the effective theory lies in
$\mathfrak B_{\rm 2t}(\kappa,q_G+q_D)$ and therefore satisfies
Theorem~\ref{thm:twotier}.
\end{theorem}

\begin{proof}
Combine \eqref{eq:BEG} and \eqref{eq:BED} by the triangle inequality and apply
the definition of the two-tier basin.
\end{proof}

This theorem is logically complete; its two hypotheses are the unresolved
four-dimensional analytic work.

\section{Three implementation branches}

\subsection{Influence branch}

Construct a block conditional transport map on the quotient by internal gauge
transformations.  On $\mathsf G_B$, derive a gauge-orthogonal Hessian or
geodesic convexity estimate and bound the derivative of the conditional center
with respect to physical boundary data.  The target is \eqref{eq:BEG}, not
mere convexity of a sampling dynamics.

\subsection{Polymer branch}

Decompose the full measure into good blocks and connected defects.  Establish
activities $z_k(X)$ satisfying, for a scale-independent $\alpha>0$,
\begin{equation}\label{eq:KP}
 \sup_B\sum_{X\ni B}|z_k(X)|\e^{\alpha|X|}\le\eta<1.
\end{equation}
Then construct $R^{(k)}$ from polymers connecting $B$ to $C$.  The key burden
is to obtain \eqref{eq:KP} along the weak-coupling trajectory and through the
finite physical entry scale, not only in the strong-coupling phase.

\subsection{Transfer branch}

Use Proposition~\ref{prop:rp} to construct the physical finite-lattice transfer
operator.  Seek a block-cone or projective contraction after quotienting gauge
null directions.  Any contraction constant must be uniform in volume and
convert to physical units before it can enter a spectral conclusion.

These branches share the same defect and reconstruction ledger.  Running them
independently provides adversarial cross-checks: a local coupling estimate that
loses positivity cannot be repaired by calling it a transfer estimate, and a
polymer estimate restricted to the trivial orbit chart does not control the
full measure.

\section{Relation to existing rigorous regimes}

Balaban's four-dimensional RG analysis supplies detailed small-field coupling
renormalization but explicitly restricts that stage to a small-field
approximation \cite{BalabanI}.  Strong-coupling lattice results can obtain
functional inequalities and exponential decay in explicit regimes
\cite{ShenZhuZhu}.  These are important calibrations on opposite sides of the
programme:
\begin{itemize}
\item the former informs the good-field multiscale construction;
\item the latter shows that local coupling and functional-inequality methods
      can yield regulated infinite-volume decay; and
\item neither result proves \eqref{eq:BEG}--\eqref{eq:BED} along the
      four-dimensional continuum trajectory.
\end{itemize}

The replacement basin is designed precisely to expose the bridge between
those regimes: regular weak-field transport plus full-measure defect
subcriticality at a finite physical scale.

\section{Frontier audit}

\begin{longtable}{p{.08\textwidth}p{.46\textwidth}p{.18\textwidth}p{.16\textwidth}}
\toprule
ID & Statement & Status & Claim level\\
\midrule
\endhead
HK1 & finite heat-kernel gauge invariance and reflection-positive factor & \Pstatus & C2 regulated\\
HK2 & fixed-translate heat-kernel TV tends to one & \Pstatus & C2\\
HK3 & naive worst-case TV basin is cutoff-uniform & \Fstatus & rejected\\
TB1 & two-tier local certificate implies a Wasserstein influence bound & \Pstatus & C2\\
TB2 & weighted two-tier basin implies clustering & \Cstatus & C2--C3\\
BE-G & four-dimensional good-field block contraction & \Ostatus & frontier\\
BE-D & four-dimensional defect/topology subcriticality & \Ostatus & frontier\\
CONT & nontrivial regulator-independent OS-positive continuum theory & \Ostatus & frontier\\
\bottomrule
\end{longtable}

The strongest advance is not a mass-gap proof.  It is a rigorous elimination
of an over-strong basin and a sharper decomposition of the first missing
nonperturbative theorem.  The new target is both more plausible and more
demanding: it must show that regular response and exceptional propagation fit
inside one contraction budget at a finite physical scale.

\section{Implemented diagnostics and immediate proof obligations}

\begin{enumerate}[label=\textbf{N\arabic*.},leftmargin=3em]
\item[\textbf{N1.}] \Pstatus\ The $SU(2)$ character-series translation
      calibration is implemented and normalized against Haar measure.
\item[\textbf{N2.}] \Pstatus\ Exact $U(1)$ convolution blocking is
      implemented.  The next step is a genuinely interacting two-plaquette or
      $2^4$ block conditional kernel on quotient boundary data.
\item[\textbf{N3.}] \Ostatus\ Test whether quotient-Wasserstein response
      decreases with interacting block size after removing pure gauge modes.
\item[\textbf{N4.}] \Pstatus\ A one-plaquette good-field tail is implemented.
      The open theorem is a connected, boundary-uniform large-deviation bound
      for \eqref{eq:goodblock} with topology recorded.
\item[\textbf{N5.}] \Pstatus\ The first finite-scale optimistic $K+R$ proxy
      is measured and fails the degree-eight threshold.  Replace the proxy by
      a certified block coupling and defect-disagreement matrix.
\item[\textbf{N6.}] \Pstatus\ A curvature-matched $U(1)$ one-plaquette
      heat-kernel/Wilson comparison is implemented.  Extend it to common
      block observables with volume- and cutoff-uniform errors.
\end{enumerate}

\begin{thebibliography}{99}
\bibitem{JaffeWitten}
A. Jaffe and E. Witten,
\emph{Quantum Yang--Mills Theory}, Clay Mathematics Institute problem
description,\newline
\url{https://www.claymath.org/wp-content/uploads/2022/06/yangmills.pdf}.

\bibitem{MenottiOnofri}
P. Menotti and E. Onofri,
``The Action of $SU(N)$ Lattice Gauge Theory in Terms of the Heat Kernel on
the Group Manifold,'' \emph{Nuclear Physics B} \textbf{190} (1981), 288--300.

\bibitem{RebeschiniVanHandel}
P. Rebeschini and R. van Handel,
``Comparison Theorems for Gibbs Measures,''
\emph{Journal of Statistical Physics} \textbf{157} (2014), 234--281,\newline
\url{https://web.math.princeton.edu/~rvan/dobrushin130819.pdf}.

\bibitem{BalabanI}
T. Ba\l aban,
``Renormalization Group Approach to Lattice Gauge Field Theories I:
Generation of Effective Actions in a Small Field Approximation and a Coupling
Constant Renormalization in Four Dimensions,''
\emph{Communications in Mathematical Physics} \textbf{109} (1987), 249--301,\newline
\url{https://projecteuclid.org/euclid.cmp/1104116842}.

\bibitem{ShenZhuZhu}
H. Shen, R. Zhu, and X. Zhu,
``A Stochastic Analysis Approach to Lattice Yang--Mills at Strong Coupling,''
\emph{Communications in Mathematical Physics} \textbf{400} (2023), 805--851,\newline
\url{https://arxiv.org/abs/2204.12737}.
\end{thebibliography}

\end{document}
